\section{Related Work}

\subsection{Language Models, Prompts, and Language Agents}

Language model is a family of machine learning model that is trained to evaluate the probability of sequences of words or tokens. Large language models (LLMs)~\citep{gpt,gpt2,gpt3,gpt3.5,gpt4,llama,llama2}  often refer to language models that adopt the autoregressive probability factorization scheme, parametrized by the Transformer architecture~\citep{vaswani2017attention}, consists of a large amount of parameters, and trained on large-scale corpus. With scaling of model size, training data, and computation, LLMs have demonstrated remarkable capabilities in generating human-like texts and understanding context. 

Prompts, on the other hand, is the key for unleashing the capabilites of LLMs. Prompts are critical components in controlling the behavior and output of LLMs and serve as the interface between human and LLMs. The design of prompts significantly impacts the performance of language models and a number of progress have been made on prompt engineering, including in-context learning~\citep{gpt3}, chain-of-thought prompting~\citep{nye2022show, wei2022chain}, ReAct~\citep{yao2022react}, self-refine~\citep{madaan2023selfrefine}, self-consistency~\citep{wang2023selfconsistency}, recurrent prompting~\citep{zhou2023recurrentgpt}, etc.

Language agents further extend the functionality of language models beyond simple prompting by allowing LLMs to use tools~\citep{schick2023toolformer} and integrating LLMs into broader systems capable of executing multi-step tasks~\citep{park2023generative,hong2023metagpt,zhou2023agents,chen2023agentverse,OpenAgents}. By stacking prompts and tools into carefully designed pipeline, agents are versatile in various applications, from customer service automation to advanced data analysis. 

\subsection{From Automated Prompt Engineering to Agent Optimization}

With the increasing popularity of prompt engineering in both academic and industry, a number of recent work investigated methods to automate the prompt engineering process. For example, \citet{pryzant2020automatically} and \citet{yang2024large} uses carefully designed prompts to unleash LLMs' ability to do prompt engineering for themselves. On the other hand, \citet{prasad-etal-2023-grips} and \citet{guo2024connecting} employs different search algorithms such as genetic algorithms for prompt optimization.

Since prompts are critical components of agents, the success of automated prompt engineering opens up the possibility of automated agent optimization. Similar to the case in automated prompt engineering, methods for agent optimization can also be categorized into two categories: \textit{prompt-based} and \textit{search-based}. For example, Agent-pro~\citep{zhang2024agentpro} and AgentOptimizer~\citep{zhang2024offline} leverage carefully designed prompts to optimize either the prompts or the tools in a node of the agent pipeline. These methods work on isolated components within an agent. Another line of research explored search-based agent optimization algorithms. ~\citet{sordoni2023joint} uses variational inference to optimize stacked LLMs. DSpy~\citep{khattab2023dspy} uses search algorithms to find the best prompts or nodes in a combinatory space. GPTSwarm~\citep{zhuge2024language} further improved the search algorithm for the combinatory optimization problem. These approaches have a few major limitations. First, the search algorithm mainly works when the metric can be defined numerically with equations that can be coded. However, most agentic tasks are real-world complex problems of which the success can not be defined by some equations, such as software development or creative writing. Second, these approaches update each component separately and therefore suffer from the local optimum of each node or component. These approaches also lack the functionality of adding nodes in the pipeline or implementing new tools.  Our proposed \baby framework, on the other hand, is the first agent learning method that optimize the agent system ``holistically'' and is able to optimize prompts, tools, nodes, as well as the way they are stacked into agents. 

Furthermore, a number of recent efforts have been done on synthesizing data to fine-tune the LLM backbone of an agent~\citep{chen2023fireact,DBLP:journals/corr/abs-2401-05268,song2024trial}. This line of research is orthogonal to our work and we believe they can be complementary to each other. ICE~\citep{qian2024investigateconsolidateexploit} is also a related work investigating inter-task transfer learning for language agents, which can be complementary with our method for building self-evolving agents.
